\annexsubsection{\texorpdfstring{Hilbert--Schmidt mappings \(^{(2)}\)}{Hilbert--Schmidt mappings (2)}}

Let E be a real Hilbert space, in which the scalar product is denoted by \((x|y)\) . There exists an isometry \(j_{E}\) of E onto its dual, characterized by the formula

\[
(x | y) = \langle x, j _ {\mathrm{E}} (y) \rangle \quad \text { for } x, y \text { in } \mathrm{E}\tag{3}
\]

(TVS, V, §1, No.~7, Th. 3).

Let \(E_{1}\) and \(E_{2}\) be two real Hilbert spaces and u a continuous linear mapping of \(E_{1}\) into \(E_{2}\) . One calls adjoint of u the continuous linear mapping \(u^{*} = j_{E_{1}}^{-1} \circ {}^{t}u \circ j_{E_{2}}\) of \(E_{2}\) into \(E_{1}\) . The mapping \(u^{*}\) is characterized by the relation

\[
(u (x _ {1}) | x _ {2}) = (x _ {1} | u ^ {*} (x _ {2})) \quad \text { for } x _ {1} \in \mathrm{E} _ {1}, x _ {2} \in \mathrm{E} _ {2}.\tag{4}
\]

If \(v\) is a continuous linear mapping of \(\mathbf{E}_2\) into a Hilbert space \(\mathbf{E}_3\) , one has \((v \circ u)^* = u^* \circ v^*\) .

Let \(E_{1}\) and \(E_{2}\) be two real Hilbert spaces and u a linear mapping of \(E_{1}\) into \(E_{2}\) . One defines two positive quadratic forms H and Q on \(E_{1}\) by the formulas

\[
\mathrm{H} (x) = \| x \| ^ {2}, \quad \mathrm{Q} (x) = \| u (x) \| ^ {2} \qquad (x \in \mathrm{E} _ {1}).
\]

PROPOSITION 3. --- Suppose u is continuous. Let \((e_{i})_{i\in I}\) be an orthonormal basis of \(E_{1}\) and \((f_{j})_{j\in J}\) an orthonormal basis of \(E_{2}\) . Then

\[
\operatorname{Tr} (\mathrm{Q} | \mathrm{H}) = \sum_ {i \in \mathrm{I}} \| u (e _ {i}) \| ^ {2} = \sum_ {j \in \mathrm{J}} \| u ^ {*} (f _ {j}) \| ^ {2} = \sum_ {i \in \mathrm{I}} \sum_ {j \in \mathrm{J}} (u (e _ {i}) | f _ {j}) ^ {2}.
\]

For every \(x \in \mathbf{E}_1\) , we have \(\|x\|^2 = \sum_{i \in \mathbf{I}} (x|e_i)^2\) , and similarly \(\|y\|^2 = \sum_{j \in \mathbf{J}} (y|f_j)^2\) for every \(y \in \mathbf{E}_2\) . Consequently,

\[
\begin{array}{r l} \sum_ {i \in \mathrm{I}} \| u (e _ {i}) \| ^ {2} & = \sum_ {i \in \mathrm{I}} \sum_ {j \in \mathrm{J}} (u (e _ {i}) | f _ {j}) ^ {2} \\ & = \sum_ {j \in \mathrm{J}} \sum_ {i \in \mathrm{I}} (e _ {i} | u ^ {*} (f _ {j})) ^ {2} \\ & = \sum_ {j \in \mathrm{J}} \| u ^ {*} (f _ {j}) \| ^ {2}. \end{array}
\]

In particular, the number \(\sum_{i\in I}\| u(e_i)\|^2\) is independent of the orthonormal basis \((e_i)_{i\in I}\) of \(\mathbf{E}_1\) .

Set \(t = \operatorname{Tr}(\mathbf{Q}|\mathbf{H})\) . For every finite subset \(I'\) of \(I\) , by definition

\[
\sum_ {i \in I ^ {\prime}} \| u (e _ {i}) \| ^ {2} = \sum_ {i \in I ^ {\prime}} Q (e _ {i}) \leqslant t,
\]

whence \(\sum_{i\in I}\| u(e_i)\|^2\leqslant t\) . Let \((e_1',\dots ,e_p')\) be a finite orthonormal sequence in \(\mathbf{E}_1\) . This sequence can be completed to an orthonormal basis \((e_{\alpha}^{\prime})_{\alpha \in A}\) of \(\mathbf{E}_1\) . Then

\[
\sum_ {\alpha = 1} ^ {p} \left\| u (e _ {\alpha} ^ {\prime}) \right\| ^ {2} \leqslant \sum_ {\alpha \in \mathrm{A}} \left\| u (e _ {\alpha} ^ {\prime}) \right\| ^ {2} = \sum_ {i \in \mathrm{I}} \left\| u (e _ {i}) \right\| ^ {2}
\]

and, passing to the supremum over all \((e_1', \ldots, e_p')\) , one obtains the inequality \(t \leqslant \sum_{i \in I} \| u(e_i) \|^2\) . We have thus established the equality \(t = \sum_{i \in I} \| u(e_i) \|^2\) .

Q.E.D.

One says that \(u\) is a Hilbert-Schmidt mapping of \(\mathbf{E}_1\) into \(\mathbf{E}_2\) if the positive quadratic form \(\mathbf{Q}: x \mapsto \| u(x) \|^2\) on \(\mathbf{E}_1\) is nuclear. When this is so, one has \(\mathbf{Q} \leqslant \operatorname{Tr}(\mathbf{Q}/\mathbf{H}) \cdot \mathbf{H}\) , therefore \(u\) is continuous and

\[
\| u \| \leqslant \operatorname{Tr} (\mathrm{Q/H}) ^ {1 / 2}.
\]

Let \(u: \mathbf{E}_1 \to \mathbf{E}_2\) be a continuous linear mapping. By Prop. 3, \(u\) is a Hilbert--Schmidt mapping if and only if there exists an orthonormal basis \((e_i)_{i \in \mathrm{I}}\) of \(\mathbf{E}_1\) such that \(\sum_{i \in \mathrm{I}} \|u(e_i)\|^2 < +\infty\) . When this is so, every orthonormal basis of \(\mathbf{E}_1\) has the same property. Moreover, if \(u\) is a Hilbert--Schmidt mapping, then so is its adjoint \(u^*\) by virtue of the formula \(\sum_{i \in \mathrm{I}} \|u(e_i)\|^2 = \sum_{j \in \mathrm{J}} \|u^*(f_j)\|^2\) (Prop. 3).

